\documentclass[10pt,a4paper]{amsart}
\usepackage{amsmath,amssymb,amsthm,graphicx}
\usepackage{verbatim}
\newenvironment{reference}{\comment}{\endcomment}
%\newenvironment{reference}{}{}
\newenvironment{slogan}{\comment}{\endcomment}
\newenvironment{history}{\comment}{\endcomment}

% For commutative diagrams we use Xy-pic
\usepackage[all]{xy}

% We use 2cell for 2-commutative diagrams.
\xyoption{2cell}
% Theorem environments.
%
\theoremstyle{plain}
\newtheorem{theorem}[subsection]{Theorem}
\newtheorem{proposition}[subsection]{Proposition}
\newtheorem{lemma}[subsection]{Lemma}

\theoremstyle{definition}
\newtheorem{definition}[subsection]{Definition}
\newtheorem{example}[subsection]{Example}
\newtheorem{exercise}[subsection]{Exercise}
\newtheorem{situation}[subsection]{Situation}

\theoremstyle{remark}
\newtheorem{remark}[subsection]{Remark}
\newtheorem{remarks}[subsection]{Remarks}

\numberwithin{equation}{subsection}

% Macros
%
\def\lim{\mathop{\mathrm{lim}}\nolimits}
\def\colim{\mathop{\mathrm{colim}}\nolimits}
\def\Spec{\mathop{\mathrm{Spec}}}
\def\Hom{\mathop{\mathrm{Hom}}\nolimits}
\def\Ext{\mathop{\mathrm{Ext}}\nolimits}
\def\SheafHom{\mathop{\mathcal{H}\!\mathit{om}}\nolimits}
\def\SheafExt{\mathop{\mathcal{E}\!\mathit{xt}}\nolimits}
\def\Sch{\mathit{Sch}}
\def\Mor{\mathop{\mathrm{Mor}}\nolimits}
\def\Ob{\mathop{\mathrm{Ob}}\nolimits}
\def\Sh{\mathop{\mathit{Sh}}\nolimits}
\def\NL{\mathop{N\!L}\nolimits}
\def\CH{\mathop{\mathrm{CH}}\nolimits}
\def\proetale{{pro\text{-}\acute{e}tale}}
\def\etale{{\acute{e}tale}}
\def\QCoh{\mathit{QCoh}}
\def\Ker{\mathop{\mathrm{Ker}}}
\def\Im{\mathop{\mathrm{Im}}}
\def\Coker{\mathop{\mathrm{Coker}}}
\def\Coim{\mathop{\mathrm{Coim}}}

% Boxtimes
%
\DeclareMathSymbol{\boxtimes}{\mathbin}{AMSa}{"02}

%
% Macros for moduli stacks/spaces
%
\def\QCohstack{\mathcal{QC}\!\mathit{oh}}
\def\Cohstack{\mathcal{C}\!\mathit{oh}}
\def\Spacesstack{\mathcal{S}\!\mathit{paces}}
\def\Quotfunctor{\mathrm{Quot}}
\def\Hilbfunctor{\mathrm{Hilb}}
\def\Curvesstack{\mathcal{C}\!\mathit{urves}}
\def\Polarizedstack{\mathcal{P}\!\mathit{olarized}}
\def\Complexesstack{\mathcal{C}\!\mathit{omplexes}}
% \Pic is the operator that assigns to X its picard group, usage \Pic(X)
% \Picardstack_{X/B} denotes the Picard stack of X over B
% \Picardfunctor_{X/B} denotes the Picard functor of X over B
\def\Pic{\mathop{\mathrm{Pic}}\nolimits}
\def\Picardstack{\mathcal{P}\!\mathit{ic}}
\def\Picardfunctor{\mathrm{Pic}}
\def\Deformationcategory{\mathcal{D}\!\mathit{ef}}
\usepackage{cleveref}
\crefname{originalreference}{original reference}{original references}
\Crefname{originalreference}{Original reference}{Original references}
\makeatletter
\newlabel{homology-definition-injective-surjective}{{010B}{source}{}{sourcepointer-1}{}}
\newlabel{homology-definition-injective-surjective@cref}{{[originalreference][0][]010B}{[1][0][]source}{}{}{}}
\makeatother
\crefname{originalreference}{original reference}{original references}
\Crefname{originalreference}{Original reference}{Original references}
\makeatletter
\newlabel{sites-definition-presheaves-injective-surjective}{{00V6}{source}{}{sourcepointer-2}{}}
\newlabel{sites-definition-presheaves-injective-surjective@cref}{{[originalreference][0][]00V6}{[1][0][]source}{}{}{}}
\makeatother
\crefname{originalreference}{original reference}{original references}
\Crefname{originalreference}{Original reference}{Original references}
\makeatletter
\newlabel{sites-definition-sheaves-injective-surjective}{{00WM}{source}{}{sourcepointer-3}{}}
\newlabel{sites-definition-sheaves-injective-surjective@cref}{{[originalreference][0][]00WM}{[1][0][]source}{}{}{}}
\makeatother
\crefname{originalreference}{original reference}{original references}
\Crefname{originalreference}{Original reference}{Original references}
\makeatletter
\newlabel{homology-lemma-adjoint-get-abelian}{{03A3}{source}{}{sourcepointer-4}{}}
\newlabel{homology-lemma-adjoint-get-abelian@cref}{{[originalreference][0][]03A3}{[1][0][]source}{}{}{}}
\makeatother
\crefname{originalreference}{original reference}{original references}
\Crefname{originalreference}{Original reference}{Original references}
\makeatletter
\newlabel{sites-section-sheaves-algebraic-structures}{{00YR}{source}{}{sourcepointer-5}{}}
\newlabel{sites-section-sheaves-algebraic-structures@cref}{{[originalreference][0][]00YR}{[1][0][]source}{}{}{}}
\makeatother
\crefname{originalsection}{original section}{original sections}
\Crefname{originalsection}{Original section}{Original sections}
\makeatletter
\newlabel{section-abelian-pre-sheaves}{{03A6}{source}{}{sourcepointer-6}{}}
\newlabel{section-abelian-pre-sheaves@cref}{{[originalsection][0][]03A6}{[1][0][]source}{}{}{}}
\makeatother
\title{The category of abelian sheaves on a site is abelian, with kernel and cokernel described by presheaf operations and sheafification}
\author{The Stacks Project Authors}
\date{stacks-project:sites-modules.tex:local-snapshot-20261012; GNU FDL 1.2 or later}
\begin{document}
\maketitle
\noindent\textit{Editorial note.} All delivered mathematical text is the author's own, copied verbatim from the pinned Stacks Project chapter file sites-modules.tex; no statement or proof step was rewritten, reordered or removed.\par\medskip
\noindent\textit{Editorial note.} The counted span cites three external results (including a section), all recorded in assembly.external\_references with their canonical tags as navigation pointers; the load-bearing reduction is the author's own verification of the criterion's hypotheses, which is inside the span.\par\medskip
\noindent\textit{Allowed standard prerequisites (editorial).}\begin{itemize}
\item Adjoint functors yielding an abelian structure on the target category. Stacks Project, Homology, Lemma homology-lemma-adjoint-get-abelian (tag 010B). The general criterion applied to the inclusion and sheafification pair, from which abelianness and the kernel/cokernel descriptions follow.
\item Sheafification of abelian presheaves: adjointness and left exactness. Stacks Project, Sites, Section sites-section-sheaves-algebraic-structures (tag 00YR). Supplies the properties of sheafification that the criterion's hypotheses require.
\end{itemize}\medskip
\ifdefined\hypertarget\hypertarget{sourcepointer-1}{}\fi
\noindent\textit{Source reference.} Original reference 010B: sources/stacks-chow/tags.txt, source lines 1321--1321. Cross-chapter standard prerequisite cited by the author inside the counted span; it is used as a pointer and its own proof is not part of this item.\par\medskip
\ifdefined\hypertarget\hypertarget{sourcepointer-2}{}\fi
\noindent\textit{Source reference.} Original reference 00V6: sources/stacks-chow/tags.txt, source lines 1136--1136. Cross-chapter standard prerequisite cited by the author inside the counted span; it is used as a pointer and its own proof is not part of this item.\par\medskip
\ifdefined\hypertarget\hypertarget{sourcepointer-3}{}\fi
\noindent\textit{Source reference.} Original reference 00WM: sources/stacks-chow/tags.txt, source lines 1188--1188. Cross-chapter standard prerequisite cited by the author inside the counted span; it is used as a pointer and its own proof is not part of this item.\par\medskip
\ifdefined\hypertarget\hypertarget{sourcepointer-4}{}\fi
\noindent\textit{Source reference.} Original reference 03A3: sources/stacks-chow/tags.txt, source lines 4294--4294. Cross-chapter standard prerequisite cited by the author inside the counted span; it is used as a pointer and its own proof is not part of this item.\par\medskip
\ifdefined\hypertarget\hypertarget{sourcepointer-5}{}\fi
\noindent\textit{Source reference.} Original reference 00YR: sources/stacks-chow/tags.txt, source lines 1265--1265. Cross-chapter standard prerequisite cited by the author inside the counted span; it is used as a pointer and its own proof is not part of this item.\par\medskip
\ifdefined\hypertarget\hypertarget{sourcepointer-6}{}\fi
\noindent\textit{Source reference.} Original section 03A6: sources/stacks-chow/tags.txt, source lines 4297--4297. Cross-chapter standard prerequisite cited by the author inside the counted span; it is used as a pointer and its own proof is not part of this item.\par\medskip
\begin{lemma}
\label{lemma-abelian-abelian}
Let $\mathcal{C}$ be a site. Let $\varphi : \mathcal{F} \to \mathcal{G}$
be a morphism of abelian sheaves on $\mathcal{C}$.
\begin{enumerate}
\item The category $\textit{Ab}(\mathcal{C})$ is an abelian category.
\item The kernel $\Ker(\varphi)$ of $\varphi$ is the same as the
kernel of $\varphi$ as a morphism of presheaves.
\item The morphism $\varphi$ is injective
(Homology, Definition \ref{homology-definition-injective-surjective})
if and only if $\varphi$ is injective as a map of presheaves
(Sites, Definition \ref{sites-definition-presheaves-injective-surjective}),
if and only if $\varphi$ is injective as a map of sheaves
(Sites, Definition \ref{sites-definition-sheaves-injective-surjective}).
\item The cokernel $\Coker(\varphi)$ of $\varphi$ is the sheafification
of the cokernel of $\varphi$ as a morphism of presheaves.
\item The morphism $\varphi$ is surjective
(Homology, Definition \ref{homology-definition-injective-surjective})
if and only if $\varphi$ is surjective as a map of sheaves
(Sites, Definition \ref{sites-definition-sheaves-injective-surjective}).
\item A complex of abelian sheaves
$$
\mathcal{F} \to \mathcal{G} \to \mathcal{H}
$$
is exact at $\mathcal{G}$ if and only if for all
$U \in \Ob(\mathcal{C})$ and all $s \in \mathcal{G}(U)$
mapping to zero in $\mathcal{H}(U)$ there exists a covering
$\{U_i \to U\}_{i \in I}$ in $\mathcal{C}$ such that each
$s|_{U_i}$ is in the image of $\mathcal{F}(U_i) \to \mathcal{G}(U_i)$.
\end{enumerate}
\end{lemma}

\begin{proof}
We claim that Homology, Lemma \ref{homology-lemma-adjoint-get-abelian}
applies to the categories $\mathcal{A} = \textit{Ab}(\mathcal{C})$
and $\mathcal{B} = \textit{PAb}(\mathcal{C})$, and the functors
$a : \mathcal{A} \to \mathcal{B}$ (inclusion),  and
$b : \mathcal{B} \to \mathcal{A}$ (sheafification).
Let us check the assumptions of
Homology, Lemma \ref{homology-lemma-adjoint-get-abelian}.
Assumption (1) is that $\mathcal{A}$, $\mathcal{B}$ are additive categories,
$a$, $b$ are additive functors, and $a$ is right adjoint to $b$.
The first two statements are clear and adjointness is
Sites, Section \ref{sites-section-sheaves-algebraic-structures} ($\epsilon$).
Assumption (2) says that $\textit{PAb}(\mathcal{C})$ is abelian
which we saw in Section \ref{section-abelian-pre-sheaves} and
that sheafification is left exact, which is
Sites, Section \ref{sites-section-sheaves-algebraic-structures} ($\zeta$).
The final assumption is that $ba \cong \text{id}_\mathcal{A}$ which is
Sites, Section \ref{sites-section-sheaves-algebraic-structures} ($\delta$).
Hence Homology, Lemma \ref{homology-lemma-adjoint-get-abelian}
applies and we conclude that $\textit{Ab}(\mathcal{C})$ is abelian.

\medskip\noindent
In the proof of Homology, Lemma \ref{homology-lemma-adjoint-get-abelian}
it is shown that $\Ker(\varphi)$ and $\Coker(\varphi)$
are equal to the sheafification of the kernel and cokernel of $\varphi$
as a morphism of abelian presheaves. This proves (4). Since the kernel
is a equalizer (i.e., a limit) and since sheafification commutes with
finite limits, we conclude that (2) holds.

\medskip\noindent
Statement (2) implies (3). Statement (4) implies (5) by our description
of sheafification. The characterization of exactness in (6) follows from
(2) and (5), and the fact that the sequence is exact if and only if
$\Im(\mathcal{F} \to \mathcal{G}) =
\Ker(\mathcal{G} \to \mathcal{H})$.
\end{proof}

\end{document}
